\chapter*{Overview}

\paragraph{The idea.} The standard way to make the hydrogen atom rigorous is
analytic: unbounded self-adjoint operators, spectral theory, all the
machinery needed to give the Coulomb Hamiltonian a precise meaning. That work
is real and hard, and Mathlib's Physlib is already making progress on it.

There is an older, more algebraic story too. The hydrogen atom has a hidden
symmetry: beyond ordinary rotations, its bound states also carry a conserved
Runge--Lenz vector, and together these generate a bigger symmetry algebra.
Lean on the representation theory of that algebra hard enough, and the
energy levels and their degeneracies fall out of group theory almost without
touching an actual Hamiltonian --- this is the classical explanation (going
back to Pauli) of why hydrogen's degeneracy is $n^2$ rather than merely
$2\ell+1$.

This blueprint tries to build that second route: work out the algebra
completely first --- classification of irreducible representations, Casimir
elements, the works --- and only afterward come back and show that the real,
concrete quantum-mechanical operators satisfy whatever the algebra needed
all along. That last step is deliberately \emph{not} attempted here (see the
note on Physlib below); the goal of this pass is to nail down the algebra and
the exact shape of the bridge to physics that will eventually be needed.

That structure splits the blueprint naturally into two halves, developed
mostly independently, that need to be stitched together at the end:
\begin{itemize}
  \item A \textbf{math half} with no physics content at all (Chapters~1--3,
    below): representations of Lie algebras, the classification of
    finite-dimensional irreducible $\sltwo$-modules, Casimir elements, and
    the classification of jointly-irreducible \emph{commuting pairs} of
    $\sltwo$-actions --- the algebraic heart of why hydrogen's degeneracy is
    $n^2$.
  \item A \textbf{physics half} that names the same objects the way a
    physicist would (Chapters~4--5, below): angular momentum operators, the
    Runge--Lenz vector, the hidden $\sofour$ symmetry of the bound states ---
    currently a self-contained finite-dimensional algebraic model,
    deliberately decoupled from Physlib's Hilbert-space machinery.
\end{itemize}

A reader with a physics background can read Chapters~2--3 as ``the abstract
reason spin is quantized in half-integer steps and total angular momentum
squared is a good quantum number,'' even though neither chapter mentions
physics until Chapter~4. A reader with a math background can read Chapters~4--5
as a worked application of Chapters~1--3 with no new mathematical content,
just new names for old objects, plus one genuinely new theorem (the
classification of commuting $\sltwo$-pairs, \S\ref{sec:sl2-commuting}) that
the application needs and that pure representation theory would want anyway.

This blueprint is organized as a dependency graph of \emph{theorems and
constructions}, not of individual small lemmas: each node below corresponds
to a genuine mathematical statement, and the graph edges (via \texttt{uses})
record exactly which results feed into which.

Nodes already available in Mathlib are marked with a dark green border in the
dependency graph (\mathlibok); nodes already proved in this project are
green (\leanok); everything else is a target for future formalization work.
Most of the blueprint is currently in this last category: the point of
publishing it at this stage is to fix the right statements and the right
dependency structure before doing the harder work of proving them, and
feedback on the definitions themselves is very welcome.
The chapters mirror the source tree:
\begin{itemize}
  \item \textbf{Representation theory} ---
    \texttt{Representation/Basic.lean}, \texttt{Representation/Schur.lean}
  \item \textbf{The Lie algebra $\sltwo$} ---
    \texttt{Sl2/Basic.lean}, \texttt{Sl2/Classification.lean},
    \texttt{Sl2/CommutingActions.lean}, \texttt{Sl2/ClebschGordan.lean}
  \item \textbf{The Casimir element} ---
    \texttt{Casimir/Basic.lean}, \texttt{Casimir/Centrality.lean}
  \item \textbf{Angular momentum} ---
    \texttt{AngularMomentum/Basic.lean}, \texttt{AngularMomentum/Classification.lean}
  \item \textbf{The hydrogen atom} ---
    \texttt{Hydrogen/Basic.lean}, \texttt{Hydrogen/Symmetry.lean},
    \texttt{Hydrogen/Spectrum.lean}
\end{itemize}

Throughout, $K$ is a field (eventually specialized to $\mathbb{C}$), $L$ a
Lie algebra over $K$, and modules are finite-dimensional unless stated
otherwise.

\paragraph{Scope of Chapters 4--5, and a note on Physlib.} Physlib already
has a substantial, mostly-proved analytic treatment of this material ---
\texttt{QuantumMechanics/Hydrogen/Basic.lean} and
\texttt{.../LaplaceRungeLenzVector.lean} define the (regularized) hydrogen
Hamiltonian and Runge--Lenz operator on Schwartz space and prove their
commutation relations and the identity $\mathbf{A}^2 = 2m\mathbf{H}(\mathbf{L}^2 +
\tfrac14\hbar^2(d-1)^2) + \dots$; \texttt{QuantumMechanics/Operators/AngularMomentum.lean}
proves the analogous facts for $\mathbf L$. Its own module documentation
lists the bound-state energy spectrum as an open TODO. Chapters~4--5 below
are deliberately kept as a \emph{decoupled, finite-dimensional algebraic
model}: they do not depend on Physlib's Hilbert-space/self-adjoint-operator
machinery, and in particular do not attempt the analytic work (essential
self-adjointness, finite-dimensionality of eigenspaces) needed to connect a
literal Physlib \texttt{HydrogenAtom} to the abstract representation theory
here. Wiring the two together --- so that this project's classification
theorems discharge Physlib's open TODO --- is future work, not attempted in
this pass.

\chapter{Representations of Lie algebras}

This chapter is pure algebra with no physics content whatsoever: it fixes
what a representation of a Lie algebra even is, and proves the single
structural fact --- Schur's lemma --- that everything downstream leans on
repeatedly. Physically-minded readers can think of a representation as ``a
symmetry acting on a space of quantum states,'' and of Schur's lemma as the
formal version of the physicist's heuristic that an operator commuting with
every symmetry of an irreducible system must be a scalar (this is exactly
how the Casimir element in Chapter~\ref{ch:casimir} is eventually shown to
act as a single number, the quantum number that labels an irreducible
multiplet).

\section{The representation object}
\label{sec:rep-basic}

Mathlib's primitive notion of a representation is not a bundled morphism but
a pair of type classes making a type $M$ into a Lie module for $L$. The
explicit homomorphism is \emph{extracted} from this structure, not the other
way around: there is currently no public constructor turning an arbitrary Lie
algebra morphism $L \to \End_K(M)$ back into these instances. Consequently
this project keeps two layers: a generic layer stated directly in terms of
\texttt{LieRingModule}/\texttt{LieModule} (used from \S\ref{sec:sl2-basic}
onward), and an explicit-morphism layer used for the physics representations
in the later chapters, connected by a bridge construction.

\begin{definition}[Lie module structure]
  \label{def:rep-lie-module}
  \lean{LieRingModule, LieModule}
  \mathlibok
  A Lie module of $L$ on a $K$-vector space $M$ is given by a bracket action
  $L \times M \to M$, $(x, m) \mapsto \liealg{x}\cdot m$, that is $K$-bilinear
  and satisfies the Leibniz identity
  $\liealg{[x,y]}\cdot m = x \cdot (y \cdot m) - y \cdot (x \cdot m)$.
\end{definition}

\begin{definition}[Representation homomorphism]
  \label{def:rep-to-end}
  \lean{LieModule.toEnd, LieModule.toEnd_apply_apply}
  \mathlibok
  \uses{def:rep-lie-module}
  Every Lie module structure on $M$ yields a Lie algebra homomorphism
  $\rho \colon L \to \End_K(M)$ with $\rho(x)(m) = \liealg{x}\cdot m$, and
  conversely the bracket action is recovered by evaluating $\rho$.
\end{definition}

\section{Generic representations as Mathlib objects}
\label{sec:rep-generic}

Mathlib does not just have the typeclass pair of
Definition~\ref{def:rep-lie-module}: it has a full bundled API around it ---
submodules, module homomorphisms, and irreducibility --- that supersedes what
a naive first pass would build from scratch. Anything stated for a type $M$
that already carries $\texttt{LieRingModule}$/$\texttt{LieModule}$ instances
(as is the case throughout Chapters~2--3, where the acting object is always
an actual Lie algebra with an actual Mathlib representation) should use
\emph{these} objects, not ad hoc predicates.

\begin{definition}[Lie submodule]
  \label{def:rep-lie-submodule}
  \lean{LieSubmodule}
  \mathlibok
  \uses{def:rep-lie-module}
  A \texttt{LieSubmodule} bundles a $K$-submodule $W \le M$ together with a
  proof that $L \cdot W \subseteq W$; this is the invariant-subspace notion,
  already packaged rather than expressed as a separate predicate.
\end{definition}

\begin{definition}[Lie module homomorphism]
  \label{def:rep-hom}
  \lean{LieModuleHom, LieModuleHom.id, LieModuleHom.comp}
  \mathlibok
  \uses{def:rep-lie-module}
  A \texttt{LieModuleHom} $M \to_{\!L} N$ (notation $M \xrightarrow{R,L} N$)
  bundles a $K$-linear map together with a proof that it commutes with the
  bracket action; it carries an identity, composition, and (since it extends
  a linear map into an $R$-module of maps) an $R$-module structure, so sums,
  differences, and scalar multiples of homomorphisms between the same pair
  of modules are again homomorphisms with no separate lemma required.
\end{definition}

\begin{theorem}[Kernel and range are Lie submodules]
  \label{thm:rep-hom-ker-range}
  \lean{LieModuleHom.ker, LieModuleHom.range}
  \mathlibok
  \uses{def:rep-hom, def:rep-lie-submodule}
  For $f \colon M \to_{\!L} N$, both $\ker f$ and $\operatorname{range} f$ are
  already \texttt{LieSubmodule}s of $M$ and $N$ respectively --- invariance
  is part of the construction, not a separate theorem.
\end{theorem}

\begin{definition}[Irreducibility]
  \label{def:rep-irreducible-generic}
  \lean{LieModule.IsIrreducible}
  \mathlibok
  \uses{def:rep-lie-submodule}
  A nontrivial Lie module $M$ is irreducible if its only \texttt{LieSubmodule}s
  are $\bot$ and $\top$ (Mathlib phrases this as $M$ being a simple order in
  the lattice of Lie submodules).
\end{definition}

\begin{definition}[Explicit representation]
  \label{def:rep-explicit}
  \lean{QuantumRepresentationTheory.LieRepresentation}
  \leanok
  \uses{def:rep-to-end}
  An explicit representation of $L$ on $V$ is a Lie algebra homomorphism
  $\rho \colon L \to \End_K(V)$. This is \emph{only} needed for the physics
  layer (Chapters~4--5), where the same underlying type must carry several
  differently-sourced actions (e.g.\ two named $\sltwo$-triples on one bound
  state space) without polluting the global instance cache with competing
  \texttt{LieModule} instances. It packages the same data as
  Definition~\ref{def:rep-to-end}, presented explicitly instead of via
  typeclass search.
\end{definition}

\begin{definition}[Invariant subspace, intertwiner (explicit form)]
  \label{def:rep-invariant}
  \lean{QuantumRepresentationTheory.IsInvariantSubspace,
        QuantumRepresentationTheory.IsIntertwiner}
  \leanok
  \uses{def:rep-explicit}
  For an explicit representation $\rho$ (Definition~\ref{def:rep-explicit}),
  a subspace $W \le V$ is invariant if $\rho(x)(W)\subseteq W$ for all
  $x \in L$, and a linear map $f\colon V\to W$ intertwines $\rho_V,\rho_W$ if
  $f\circ\rho_V(x) = \rho_W(x)\circ f$ for all $x$. These are the explicit-form
  counterparts of Definitions~\ref{def:rep-lie-submodule},
  \ref{def:rep-hom}, needed only because $\rho$ does not come with ambient
  typeclass instances.
\end{definition}

\begin{definition}[Irreducibility (explicit form)]
  \label{def:rep-irreducible}
  \lean{QuantumRepresentationTheory.IsIrreducible}
  \leanok
  \uses{def:rep-invariant}
  $\rho$ is irreducible if its only invariant subspaces are $0$ and $V$.
\end{definition}

\begin{definition}[Intertwiner (explicit form)]
  \label{def:rep-intertwiner}
  \lean{QuantumRepresentationTheory.IsIntertwiner}
  \leanok
  \uses{def:rep-explicit}
  See Definition~\ref{def:rep-invariant}.
\end{definition}

\begin{theorem}[Kernel of an intertwiner is invariant (explicit form)]
  \label{thm:rep-ker-invariant}
  \lean{QuantumRepresentationTheory.IsIntertwiner.ker_isInvariant}
  \leanok
  \uses{def:rep-intertwiner, def:rep-invariant}
  If $f$ intertwines $\rho_V$ and $\rho_W$, then $\ker f$ is a $\rho_V$-invariant
  subspace of $V$. (This is exactly Theorem~\ref{thm:rep-hom-ker-range}, restated
  for the explicit form where $\ker f$ is not automatically bundled as a
  submodule.)
\end{theorem}

\begin{theorem}[Range of an intertwiner is invariant (explicit form)]
  \label{thm:rep-range-invariant}
  \lean{QuantumRepresentationTheory.IsIntertwiner.range_isInvariant}
  \leanok
  \uses{def:rep-intertwiner, def:rep-invariant}
  If $f$ intertwines $\rho_V$ and $\rho_W$, then $\operatorname{range} f$ is a
  $\rho_W$-invariant subspace of $W$.
\end{theorem}

\begin{remark}[Missing bridge: from a morphism to a module instance]
  \label{def:rep-bridge}
  \uses{def:rep-explicit, def:rep-to-end}
  A genuinely missing piece of Mathlib API is a construction turning an
  explicit $\rho \colon L \to \End_K(V)$ (Definition~\ref{def:rep-explicit})
  into local \texttt{LieRingModule}/\texttt{LieModule} instances on $V$
  (Definition~\ref{def:rep-lie-module}), so that the generic machinery of
  \S\ref{sec:rep-generic}--\S\ref{sec:cas-centrality} can be applied to the
  explicit physics representations of Chapters~4--5. This should be phrased
  either as a \texttt{local instance}-producing definition or as a
  \texttt{LieModule.toEnd}-compatible wrapper type.
\end{remark}

\section{Schur's lemma}
\label{sec:rep-schur}

We take the direct eigenvalue route to Schur's lemma rather than routing
through the universal enveloping algebra: it is shorter, and it isolates
exactly the facts needed later (\S\ref{sec:cas-centrality}) to show the
Casimir element acts by a scalar. It is stated generically
(Definitions~\ref{def:rep-hom}, \ref{def:rep-irreducible-generic}), which is
the form Chapters~2--3 actually need.

\begin{theorem}[Existence of an eigenvalue]
  \label{thm:rep-eigenvalue-exists}
  \lean{Module.End.exists_eigenvalue}
  \mathlibok
  If $K$ is algebraically closed, $V$ is finite-dimensional and nonzero, and
  $f \in \End_K(V)$, then $f$ has an eigenvalue.
\end{theorem}

\begin{theorem}[Eigenspace of an endomorphism is nonzero]
  \label{thm:rep-eigenspace-nonzero}
  \lean{QuantumRepresentationTheory.exists_ker_smul_id_ne_bot}
  \uses{thm:rep-eigenvalue-exists}
  Under the hypotheses of Theorem~\ref{thm:rep-eigenvalue-exists}, there is
  $c \in K$ with $\ker(f - c \cdot \mathrm{id}_V) \neq 0$.
\end{theorem}

\begin{theorem}[Schur's lemma]
  \label{thm:rep-schur}
  \lean{QuantumRepresentationTheory.schur}
  \uses{def:rep-hom, def:rep-irreducible-generic, thm:rep-hom-ker-range,
        thm:rep-eigenspace-nonzero}
  Let $K$ be algebraically closed, let $M$ be a finite-dimensional
  irreducible Lie module for $L$ (Definition~\ref{def:rep-irreducible-generic}),
  and let $f \colon M \to_{\!L} M$ be a \texttt{LieModuleHom}
  (Definition~\ref{def:rep-hom}). Then $f = c \cdot \mathrm{id}_M$ for some
  $c \in K$: apply Theorem~\ref{thm:rep-eigenspace-nonzero} to the underlying
  endomorphism to get $c$ with $\ker(f - c\cdot\mathrm{id}) \neq 0$; since
  $f - c\cdot\mathrm{id}$ is again a \texttt{LieModuleHom}, its kernel is a
  Lie submodule (Theorem~\ref{thm:rep-hom-ker-range}), hence all of $M$ by
  irreducibility.
\end{theorem}

\begin{remark}[Parallel: Schur via simple modules over the enveloping algebra]
  \label{rem:rep-schur-uea}
  \uses{thm:rep-schur, def:cas-uea-rep}
  Mathlib already gives a division-ring structure on $\End_R(M)$ for any
  simple module $M$ (\lean{Module.End.instDivisionRing}\mathlibok). Once
  Definition~\ref{def:cas-uea-rep} equips $V$ with a module structure over
  the universal enveloping algebra $\UEA{L}$, and Lie-irreducibility of $V$
  is shown equivalent to simplicity as a $\UEA{L}$-module, this instance
  reproves Theorem~\ref{thm:rep-schur} as a special case. This is recorded as
  a generalization branch, not a prerequisite for the first proof of Schur's
  lemma.
\end{remark}

\chapter{The Lie algebra $\sltwo$ and its representations}
\label{ch:sl2}

$\sltwo$ is the smallest non-abelian Lie algebra, and its finite-dimensional
representations are completely classified: there is exactly one irreducible
representation of each dimension. That single fact, proved with no physics
in sight, is secretly the entire reason angular momentum is quantized in
half-integer steps: $\sltwo$ is (isomorphic to the complexification of) the
angular-momentum algebra $\liealg{so}_3$, and Chapter~\ref{ch:am} does
nothing more than transport this classification across that isomorphism. The
section on commuting $\sltwo$-actions (\S\ref{sec:sl2-commuting}) is the one
piece of genuinely new mathematics this blueprint needs and Mathlib does not
yet have: it is what turns ``two commuting conserved quantities'' into a
concrete dimension count, and it is the algebraic engine behind the
hydrogen atom's $n^2$ degeneracy in Chapter~\ref{ch:hydrogen}.

\section{$\sltwo$-triples and primitive vectors}
\label{sec:sl2-basic}

Mathlib already contains essentially the whole primitive-vector string
calculus for $\sltwo$-triples; the classification work in
\S\ref{sec:sl2-classification} only needs to assemble these pieces.

\begin{definition}[$\sltwo$-triple]
  \label{def:sl2-triple}
  \lean{IsSl2Triple}
  \mathlibok
  \uses{def:rep-lie-module}
  A triple $(h,e,f)$ of elements of $L$ is an $\sltwo$-triple if $h \neq 0$,
  $[e,f] = h$, $[h,e] = 2e$, and $[h,f] = -2f$; equivalently, the subalgebra
  they generate is a copy of $\sltwo$.
\end{definition}

\begin{definition}[Primitive vector]
  \label{def:sl2-primitive}
  \lean{IsSl2Triple.HasPrimitiveVectorWith}
  \mathlibok
  \uses{def:sl2-triple}
  Given an $\sltwo$-triple $t = (h,e,f)$ acting on a module $M$, a nonzero
  vector $m \in M$ is primitive of weight $\mu \in K$ for $t$ if
  $h \cdot m = \mu m$ and $e \cdot m = 0$.
\end{definition}

\begin{theorem}[$h$-eigenvalues along the lowering string]
  \label{thm:sl2-string-h}
  \lean{IsSl2Triple.HasPrimitiveVectorWith.lie_h_pow_toEnd_f}
  \mathlibok
  \uses{def:sl2-primitive}
  For a primitive vector $m$ of weight $\mu$, $h \cdot (f^k m) = (\mu - 2k)(f^k m)$
  for every $k \ge 0$.
\end{theorem}

\begin{theorem}[Raising along the lowering string]
  \label{thm:sl2-string-e}
  \lean{IsSl2Triple.HasPrimitiveVectorWith.lie_e_pow_succ_toEnd_f}
  \mathlibok
  \uses{def:sl2-primitive}
  For a primitive vector $m$ of weight $\mu$,
  $e \cdot (f^{k+1} m) = (k+1)(\mu - k)(f^k m)$ for every $k \ge 0$.
\end{theorem}

\begin{theorem}[The string is finite with integral weight]
  \label{thm:sl2-string-finite}
  \lean{IsSl2Triple.HasPrimitiveVectorWith.exists_nat,
        IsSl2Triple.HasPrimitiveVectorWith.pow_toEnd_f_ne_zero_of_eq_nat,
        IsSl2Triple.HasPrimitiveVectorWith.pow_toEnd_f_eq_zero_of_eq_nat}
  \mathlibok
  \uses{thm:sl2-string-h, thm:sl2-string-e}
  If $M$ is finite-dimensional (Noetherian, torsion-free, over a
  characteristic-zero domain), the weight $\mu$ of a primitive vector $m$
  equals some $n \in \mathbb{N}$, $f^k m \neq 0$ for all $k \le n$, and
  $f^{n+1} m = 0$.
\end{theorem}

\begin{theorem}[Finite-dimensional modules over an algebraically closed field are triangularizable]
  \label{thm:sl2-triangularizable}
  \lean{LieModule.instIsTriangularizableOfIsAlgClosed}
  \mathlibok
  \uses{def:rep-lie-module}
  If $K$ is algebraically closed and $M$ is a finite-dimensional Lie module
  for $L$, then $M$ is automatically triangularizable: this is an instance in
  Mathlib and requires no further formalization work for the $K = \mathbb{C}$
  case used throughout this project.
\end{theorem}

\begin{theorem}[Existence of a primitive vector]
  \label{thm:sl2-primitive-exists}
  \lean{IsSl2Triple.exists_hasPrimitiveVectorWith}
  \mathlibok
  \uses{def:sl2-triple, thm:sl2-triangularizable}
  If $K$ has characteristic zero, $M$ is a nonzero finite-dimensional,
  torsion-free, triangularizable module for an $\sltwo$-triple $t$, then $M$
  contains a primitive vector for $t$.
\end{theorem}

\begin{definition}[Standard module $V(n)$]
  \label{def:sl2-standard-module}
  \lean{QuantumRepresentationTheory.StandardSl2Module, QuantumRepresentationTheory.StandardSl2Module.v, QuantumRepresentationTheory.standardSl2ModuleLieRingModule, QuantumRepresentationTheory.standardSl2ModuleLieModule}
  For $n \in \mathbb{N}$, the standard module $V(n)$ is the
  $(n+1)$-dimensional $\sltwo$-module with basis $v_0, \dots, v_n$ on which
  $h \cdot v_k = (n-2k) v_k$, $f \cdot v_k = v_{k+1}$ (and $v_{n+1} := 0$), and
  $e \cdot v_k = k(n-k+1) v_{k-1}$. This is the explicit model produced by
  the classification theorem below.
\end{definition}

\section{Classification of irreducible $\sltwo$-modules}
\label{sec:sl2-classification}

\begin{theorem}[Linear independence of the primitive-vector string]
  \label{thm:sl2-string-independent}
  \lean{QuantumRepresentationTheory.string_linearIndependent}
  \uses{thm:sl2-string-finite, thm:sl2-string-h}
  If $m$ is primitive of weight $n \in \mathbb{N}$, the vectors
  $m, fm, \dots, f^n m$ are linearly independent, since they lie in distinct
  $h$-eigenspaces (Theorem~\ref{thm:sl2-string-h}) and are all nonzero
  (Theorem~\ref{thm:sl2-string-finite}).
\end{theorem}

\begin{theorem}[The string spans a submodule]
  \label{thm:sl2-string-submodule}
  \lean{QuantumRepresentationTheory.exists_stringLieSubmodule}
  \uses{thm:sl2-string-finite, thm:sl2-string-h, thm:sl2-string-e}
  The span of $m, fm, \dots, f^n m$ is invariant under $h$, $e$, and $f$,
  hence is a Lie submodule.
\end{theorem}

\begin{theorem}[Irreducibility forces the string to span]
  \label{thm:sl2-irreducible-spans}
  \lean{QuantumRepresentationTheory.stringSpan_eq_top}
  \uses{def:rep-irreducible-generic, thm:sl2-string-submodule}
  If $M$ is an irreducible $\sltwo$-module, the submodule of
  Theorem~\ref{thm:sl2-string-submodule} is nonzero, hence equal to $M$.
\end{theorem}

\begin{theorem}[Basis theorem]
  \label{thm:sl2-basis-theorem}
  \lean{QuantumRepresentationTheory.string_isBasis, QuantumRepresentationTheory.finrank_eq_succ}
  \uses{thm:sl2-string-independent, thm:sl2-irreducible-spans}
  If $M$ is a finite-dimensional irreducible $\sltwo$-module with primitive
  vector $m$ of weight $n$, then $m, fm, \dots, f^n m$ is a basis of $M$; in
  particular $\dim M = n+1$.
\end{theorem}

\begin{theorem}[Classification of irreducible $\sltwo$-modules]
  \label{thm:sl2-classification}
  \lean{QuantumRepresentationTheory.classification}
  \uses{thm:sl2-basis-theorem, def:sl2-standard-module,
        thm:sl2-primitive-exists}
  Every finite-dimensional irreducible $\sltwo$-module over an algebraically
  closed field of characteristic zero is isomorphic, as an $\sltwo$-module,
  to the standard module $V(n)$ (Definition~\ref{def:sl2-standard-module})
  for a unique $n \in \mathbb{N}$.
\end{theorem}

\begin{theorem}[Uniqueness of the highest weight]
  \label{thm:sl2-uniqueness-hw}
  \lean{QuantumRepresentationTheory.classification_n_eq_finrank_sub_one}
  \uses{thm:sl2-classification}
  The integer $n$ in Theorem~\ref{thm:sl2-classification} is an isomorphism
  invariant of $M$: it equals $\dim M - 1$.
\end{theorem}

\section{Commuting $\sltwo$-actions and tensor products}
\label{sec:sl2-commuting}

This is the largest hidden node in the whole roadmap. Without it, knowing
that two representations of two \emph{commuting} $\sltwo$-triples have equal
Casimir eigenvalues does not by itself give a dimension formula; it is the
key input to the hydrogen spectrum in Chapter~5, and is more fundamental to
that goal than the classical Clebsch--Gordan decomposition of
\S\ref{sec:sl2-cg}.

\begin{theorem}[Diagonal $\sltwo$-action on a tensor product]
  \label{thm:sl2-tensor-action}
  \lean{TensorProduct.LieModule.lieRingModule, TensorProduct.LieModule.lieModule}
  \mathlibok
  \uses{def:rep-lie-module}
  For Lie modules $M$, $N$ of $L$, the tensor product $M \otimes_K N$ carries
  the Lie module structure $x \cdot (m \otimes n) = (x\cdot m)\otimes n +
  m \otimes (x \cdot n)$.
\end{theorem}

\begin{theorem}[Classification of jointly irreducible commuting $\sltwo$-actions]
  \label{thm:sl2-commuting-classification}
  \lean{QuantumRepresentationTheory.commuting_classification, QuantumRepresentationTheory.commuting_classification_finrank}
  \uses{thm:sl2-classification, thm:sl2-tensor-action}
  Let $\rho_1, \rho_2$ be representations of two $\sltwo$-triples on a
  finite-dimensional space $V$ over an algebraically closed field of
  characteristic zero, with $[\rho_1(x), \rho_2(y)] = 0$ for all $x, y$, and
  suppose $V$ has no proper subspace invariant under both actions
  simultaneously. Then there exist $m, n \in \mathbb{N}$ and a linear
  isomorphism
  \[
    V \;\cong_K\; V(m) \otimes_K V(n)
  \]
  intertwining $(\rho_1, \rho_2)$ with the actions of the two tensor factors
  of Theorem~\ref{thm:sl2-tensor-action}.
\end{theorem}

\section{Clebsch--Gordan decomposition}
\label{sec:sl2-cg}

\begin{theorem}[Clebsch--Gordan decomposition]
  \label{thm:sl2-clebsch-gordan}
  \lean{QuantumRepresentationTheory.clebsch_gordan_finrank, QuantumRepresentationTheory.clebsch_gordan}
  \uses{thm:sl2-classification, thm:sl2-tensor-action}
  For $m, n \in \mathbb{N}$, the diagonal $\sltwo$-module $V(m) \otimes_K
  V(n)$ (Theorem~\ref{thm:sl2-tensor-action}) decomposes as
  \[
    V(m) \otimes_K V(n) \;\cong_K\; \bigoplus_{k=0}^{\min(m,n)} V(m+n-2k).
  \]
\end{theorem}

\chapter{The Casimir element}
\label{ch:casimir}

Physically, ``total angular momentum squared,'' $J^2 = J_x^2+J_y^2+J_z^2$, is
special because it commutes with every rotation generator, so it can be
diagonalized simultaneously with a full set of symmetry generators and used
to label irreducible multiplets by a single quantum number $j(j+1)$. The
Casimir element is the algebraic definition making this precise for an
arbitrary Lie algebra with an invariant form, not just $\liealg{so}_3$: an
element built from the generators themselves that is guaranteed to commute
with all of them, and which Schur's lemma (Chapter~1) then forces to act as
a scalar on any irreducible representation. Chapter~\ref{ch:am} identifies
$J^2$ with the $\sltwo$ Casimir element built here, and
Chapter~\ref{ch:hydrogen} uses \emph{two} Casimir elements at once ---
one for each of the two commuting $\sltwo$-actions hiding inside the
hydrogen atom's dynamical symmetry --- to pin down its spectrum.

\section{Construction}
\label{sec:cas-basic}

\begin{definition}[Nondegenerate invariant bilinear form]
  \label{def:cas-invariant-form}
  \lean{LinearMap.BilinForm.lieInvariant, LinearMap.BilinForm.Nondegenerate}
  \mathlibok
  \uses{def:rep-lie-module}
  A bilinear form $\Phi$ on $L$ is Lie-invariant if $\Phi([x,y],z) +
  \Phi(y,[x,z]) = 0$ for all $x,y,z \in L$; it is nondegenerate if its
  radical is trivial. For $L$ semisimple (over a characteristic-zero field,
  e.g.\ $\mathbb{C}$), the Killing form is such a form: Mathlib has both
  Cartan criteria in \texttt{Mathlib.Algebra.Lie.CartanCriterion}, and in
  particular \texttt{LieAlgebra.hasTrivialRadical\_iff\_isKilling} gives the
  equivalence between trivial radical (which coincides with semisimplicity
  in characteristic zero) and a nondegenerate Killing form
  (\texttt{LieAlgebra.IsKilling}).
\end{definition}

\begin{definition}[Dual basis]
  \label{def:cas-dual-basis}
  \lean{LinearMap.BilinForm.dualBasis}
  \mathlibok
  \uses{def:cas-invariant-form}
  Given a basis $x_1, \dots, x_d$ of $L$ and a nondegenerate form $\Phi$, the
  dual basis $x^1, \dots, x^d$ is the unique basis with
  $\Phi(x_i, x^j) = \delta_i^j$. (Standard linear-algebra fact about
  nondegenerate bilinear forms, available generically in Mathlib.)
\end{definition}

\begin{definition}[Basis-dependent Casimir element]
  \label{def:cas-basis-dependent}
  \lean{QuantumRepresentationTheory.casimirBasisDependent}
  \uses{def:cas-dual-basis, def:rep-to-end}
  For a representation $\rho$ of $L$ on $V$ and a basis $x_1,\dots,x_d$ of
  $L$ with $\Phi$-dual basis $x^1,\dots,x^d$, the Casimir operator is
  \[
    C_\rho \;=\; \sum_{i=1}^d \rho(x_i)\,\rho(x^i) \;\in\; \End_K(V).
  \]
\end{definition}

\begin{theorem}[Change-of-basis invariance]
  \label{thm:cas-change-of-basis}
  \lean{QuantumRepresentationTheory.casimirBasisDependent_basis_indep}
  \uses{def:cas-basis-dependent}
  $C_\rho$, as defined in Definition~\ref{def:cas-basis-dependent}, does not
  depend on the choice of basis $x_1,\dots,x_d$ of $L$.
\end{theorem}

\begin{definition}[The Casimir operator]
  \label{def:cas-casimir}
  \lean{QuantumRepresentationTheory.casimir}
  \uses{thm:cas-change-of-basis}
  Write $C_\rho \in \End_K(V)$ for the common value of
  Definition~\ref{def:cas-basis-dependent}, independent of the choice of
  basis by Theorem~\ref{thm:cas-change-of-basis}.
\end{definition}

\section{Centrality and the scalar action}
\label{sec:cas-centrality}

\begin{theorem}[The Casimir element commutes with the generators]
  \label{thm:cas-commutes-generators}
  \lean{QuantumRepresentationTheory.casimir_commutes}
  \uses{def:cas-casimir, def:cas-invariant-form}
  For every $y \in L$, $[\rho(y), C_\rho] = 0$ in $\End_K(V)$.
\end{theorem}

\begin{definition}[Universal enveloping algebra representation]
  \label{def:cas-uea-rep}
  \lean{UniversalEnvelopingAlgebra.ι, UniversalEnvelopingAlgebra.lift,
        UniversalEnvelopingAlgebra.lift_unique,
        UniversalEnvelopingAlgebra.hom_ext}
  \mathlibok
  \uses{def:rep-to-end}
  The universal enveloping algebra $\UEA{L}$, together with
  $\iota \colon L \to \UEA{L}$, has the universal property that every
  representation $\rho \colon L \to \End_K(V)$ extends uniquely to an
  associative-algebra homomorphism $\tilde\rho \colon \UEA{L} \to \End_K(V)$
  with $\tilde\rho \circ \iota = \rho$.
\end{definition}

\begin{theorem}[Centrality of the Casimir element]
  \label{thm:cas-central}
  \lean{QuantumRepresentationTheory.casimirUEA, QuantumRepresentationTheory.casimirUEA_mem_center}
  \lean{Subalgebra.center}
  \uses{thm:cas-commutes-generators, def:cas-uea-rep}
  The element $C_\rho$, viewed inside $\UEA{L}$ via $\iota$ and
  Definition~\ref{def:cas-basis-dependent} applied to the adjoint
  representation, lies in the center \texttt{Subalgebra.center K (UEA K L)}.
  (Note: the correct ambient object is \texttt{Subalgebra.center}, not
  \texttt{Subring.center}, since we need $K$-linearity of the centralizing
  condition.)
\end{theorem}

\begin{theorem}[A central element gives an intertwiner]
  \label{thm:cas-central-intertwines}
  \lean{QuantumRepresentationTheory.central_commutes_toEnd}
  \uses{thm:cas-central, def:cas-uea-rep, def:rep-hom}
  If $z \in \UEA{L}$ is central and $\tilde\rho$ is the extension of a
  representation $\rho$ on $V$ (Definition~\ref{def:cas-uea-rep}), then
  $\tilde\rho(z) \colon V \to V$, viewed as a self-map of $V$, is a
  \texttt{LieModuleHom} (Definition~\ref{def:rep-hom}): it commutes with
  $\rho(x)$ for every $x \in L$ since $z$ commutes with $\iota(x)$ in
  $\UEA{L}$.
\end{theorem}

\begin{theorem}[Schur gives a scalar Casimir]
  \label{thm:cas-scalar}
  \lean{QuantumRepresentationTheory.casimir_scalar}
  \uses{thm:rep-schur, thm:cas-central-intertwines, def:cas-casimir}
  If $\rho$ is a finite-dimensional irreducible representation of $L$ over an
  algebraically closed field, $C_\rho$ acts on $V$ as multiplication by a
  scalar $c(\rho) \in K$.
\end{theorem}

\chapter{Angular momentum}
\label{ch:am}

Physics starts here: this chapter is where Chapters~1--3's purely algebraic
results get their first physical names. The three angular-momentum operators
$J_x,J_y,J_z$ and their commutation relations are exactly an $\sltwo$-triple
in disguise, so the classification of irreducible $\sltwo$-modules
(\S\ref{sec:sl2-classification}) immediately gives the spin-quantization rule
--- every irreducible multiplet of angular-momentum states has a half-integer
spin $j$ and dimension $2j+1$ --- as a corollary, not a new proof. As
discussed in the overview, this chapter is deliberately kept as a
finite-dimensional algebraic model: the operators $J_x,J_y,J_z$ are given as
abstract data satisfying the right commutation relations, not derived from
Physlib's Hilbert-space momentum/position operators. This chapter's results
are used twice in Chapter~\ref{ch:hydrogen}: once directly (ordinary orbital
angular momentum) and once again after being ``split in half'' by the
Runge--Lenz vector into the two commuting copies that generate hydrogen's
hidden $\sofour$ symmetry.

\section{The algebraic representation}
\label{sec:am-basic}

\begin{definition}[Angular momentum representation]
  \label{def:am-representation}
  \lean{QuantumRepresentationTheory.AngularMomentumRepresentation}
  \uses{def:rep-explicit}
  An algebraic angular momentum representation is an explicit representation
  $\rho$ (Definition~\ref{def:rep-explicit}) of the Lie algebra $\liealg{so}_3
  \cong \sltwo$ on a state space $V$, presented via operators $J_x, J_y, J_z$
  satisfying $[J_x,J_y] = iJ_z$ and cyclic permutations. The physical
  requirement that $J_x,J_y,J_z$ be self-adjoint on a Hilbert space is
  \emph{not} part of this definition: it is exactly the kind of fact that
  belongs on the Physlib side of the future bridge discussed in the Overview,
  not something assumed here.

  \emph{Caveat:} the commutation relations here are currently axiomatized as
  three fields of a structure. Mathlib already packages $\liealg{so}_3$
  concretely as \texttt{Cross.lieRing}, the Lie ring structure on
  $\texttt{Fin 3} \to R$ given by the cross product (in
  \texttt{Mathlib.LinearAlgebra.CrossProduct}), with the Jacobi identity and
  antisymmetry already proved (\texttt{jacobi\_cross}, \texttt{cross\_anticomm}).
  Presenting angular momentum as an honest representation of that concrete
  Lie ring, rather than an axiomatized triple of operators, is a more
  idiomatic route not yet taken --- see the caveats chapter.
\end{definition}

\begin{definition}[Ladder operators]
  \label{def:am-ladder}
  \lean{QuantumRepresentationTheory.AngularMomentumRepresentation.Jp, QuantumRepresentationTheory.AngularMomentumRepresentation.Jm}
  \uses{def:am-representation}
  $J_\pm := J_x \pm i J_y$.
\end{definition}

\begin{theorem}[Commutation relations for the ladder operators]
  \label{thm:am-commutators}
  \lean{QuantumRepresentationTheory.AngularMomentumRepresentation.comm_z_p, QuantumRepresentationTheory.AngularMomentumRepresentation.comm_z_m, QuantumRepresentationTheory.AngularMomentumRepresentation.comm_p_m}
  \uses{def:am-ladder}
  $[J_z, J_\pm] = \pm J_\pm$ and $[J_+, J_-] = 2J_z$.
\end{theorem}

\begin{theorem}[Angular momentum generates an $\sltwo$-triple]
  \label{thm:am-is-sl2-triple}
  \lean{QuantumRepresentationTheory.AngularMomentumRepresentation.isSl2Triple}
  \uses{thm:am-commutators, def:sl2-triple, def:rep-bridge}
  $(J_z, J_+, J_-)$ is an $\sltwo$-triple in the sense of
  Definition~\ref{def:sl2-triple}, via the bridge construction of
  Remark~\ref{def:rep-bridge} turning $\rho$ into local Lie-module instances.
\end{theorem}

\section{Spin classification}
\label{sec:am-classification}

\begin{theorem}[Irreducibility transfers along the bridge]
  \label{thm:am-irreducibility-transfer}
  \lean{QuantumRepresentationTheory.AngularMomentumRepresentation.isIrreducible_iff}
  \uses{def:rep-irreducible, thm:am-is-sl2-triple}
  $\rho$ is irreducible as an explicit representation
  (Definition~\ref{def:rep-irreducible}) iff the induced $\sltwo$-module
  structure of Theorem~\ref{thm:am-is-sl2-triple} is irreducible.
\end{theorem}

\begin{theorem}[Spin classification]
  \label{thm:am-spin-classification}
  \lean{QuantumRepresentationTheory.AngularMomentumRepresentation.spin_classification}
  \uses{thm:am-irreducibility-transfer, thm:sl2-classification}
  Every finite-dimensional irreducible angular momentum representation is
  isomorphic to $V(2j)$ for a unique half-integer spin
  $j \in \tfrac12\mathbb{N}$, of dimension $2j+1$.
\end{theorem}

\begin{theorem}[$J^2$ is the Casimir element]
  \label{thm:am-j2-casimir}
  \lean{QuantumRepresentationTheory.AngularMomentumRepresentation.Jsq_eq_smul}
  \uses{def:cas-casimir, thm:am-is-sl2-triple}
  $J^2 := J_x^2+J_y^2+J_z^2$ agrees, up to the normalization of
  Definition~\ref{def:cas-basis-dependent} for $\sltwo$, with the Casimir
  operator $C_\rho$; on the spin-$j$ representation it acts as $j(j+1)$.
\end{theorem}

\begin{definition}[Magnetic-weight basis]
  \label{def:am-magnetic-basis}
  \lean{QuantumRepresentationTheory.AngularMomentumRepresentation.magneticBasisVec}
  \uses{thm:sl2-basis-theorem, thm:am-spin-classification}
  The basis $\ket{j,m}$ for $m = -j,\dots,j$ of the spin-$j$ representation
  produced by Theorem~\ref{thm:sl2-basis-theorem}, reindexed by the magnetic
  quantum number $m = j - k$.
\end{definition}

\chapter{The hydrogen atom}
\label{ch:hydrogen}

This is the payoff chapter, and where everything above gets stitched
together. Besides ordinary angular momentum $\vec J$, a classical particle
orbiting under an inverse-square force conserves a second vector quantity,
the Runge--Lenz vector, which points along the orbit's major axis; quantum
mechanically it survives as an operator $\vec A$ on the bound-state
eigenspaces. Together $\vec J$ and (a rescaled) $\vec A$ generate a bigger
symmetry algebra than rotations alone, and --- remarkably --- it splits as
two \emph{commuting} copies of the angular-momentum algebra from
Chapter~\ref{ch:am}, $\vec I = \tfrac12(\vec J+\vec A)$ and
$\vec K = \tfrac12(\vec J-\vec A)$. Feeding this pair into the classification
of commuting $\sltwo$-actions (\S\ref{sec:sl2-commuting}) and the two Casimir
elements it carries (Chapter~\ref{ch:casimir}) is what forces the two
copies to have equal spin and produces, with no reference to the radial
Schr\"odinger equation, the characteristic $n^2$-fold degenerate,
$-1/n^2$-spaced hydrogen spectrum. As throughout, this chapter is a
self-contained algebraic model: it takes the existence of a
finite-dimensional bound-state eigenspace with the right operators on it as
given, rather than deriving it from Physlib's Hamiltonian (see the note on
Physlib in the Overview).

\section{The Runge--Lenz vector}
\label{sec:hyd-basic}

\begin{definition}[Levi-Civita symbol]
  \label{def:hyd-levi-civita}
  \lean{QuantumRepresentationTheory.leviCivita}
  $\varepsilon_{ijk}$, the totally antisymmetric symbol on $\{1,2,3\}$ with
  $\varepsilon_{123}=1$, used to write the cross product and the commutation
  relations below.

  \emph{Caveat:} Mathlib already has a bilinear cross product on
  \texttt{Fin 3 → R} (\texttt{crossProduct}, in
  \texttt{Mathlib.LinearAlgebra.CrossProduct}), with bilinearity built into
  its type (it is a \texttt{LinearMap.mk₂}) rather than proved separately.
  Reusing it in place of this from-scratch Levi-Civita definition, so that
  Theorem~\ref{thm:hyd-cross-linearity} below becomes free, is a natural
  simplification not yet made --- see the caveats chapter.
\end{definition}

\begin{theorem}[Linearity of the cross-product components]
  \label{thm:hyd-cross-linearity}
  \lean{QuantumRepresentationTheory.crossComponent, QuantumRepresentationTheory.crossComponent_bilinear}
  \uses{def:hyd-levi-civita}
  Each component $(A \times B)_i = \sum_{jk} \varepsilon_{ijk} A_j B_k$ is
  bilinear in $A, B$; used repeatedly to expand commutators involving the
  Runge--Lenz vector.
\end{theorem}

\begin{definition}[Bound-state representation]
  \label{def:hyd-bound-state}
  \lean{QuantumRepresentationTheory.BoundStateRepresentation}
  \uses{def:am-representation}
  The explicit representation of the hydrogen Hamiltonian's dynamical
  symmetry algebra on a fixed negative-energy eigenspace $V_E$ of the
  Hamiltonian $H$, carrying the angular momentum representation
  $(J_x,J_y,J_z)$ of Definition~\ref{def:am-representation} together with
  the (unnormalized) Runge--Lenz operator $\vec M$.
\end{definition}

\begin{theorem}[The rescaling constant is nonzero]
  \label{thm:hyd-rescaling-nonzero}
  \lean{QuantumRepresentationTheory.BoundStateRepresentation.rescaling_pos, QuantumRepresentationTheory.BoundStateRepresentation.rescalingConst, QuantumRepresentationTheory.BoundStateRepresentation.rescalingConst_pos}
  \uses{def:hyd-bound-state}
  On $V_E$ with $E < 0$, the scalar $-2E/m$ appearing in the normalization of
  $\vec M$ is strictly positive, so
  $\vec A := \vec M / \sqrt{-2mE}$ is well defined.
\end{theorem}

\begin{definition}[Rescaled Runge--Lenz vector]
  \label{def:hyd-rescaled-rl}
  \lean{QuantumRepresentationTheory.BoundStateRepresentation.Ax, QuantumRepresentationTheory.BoundStateRepresentation.Ay, QuantumRepresentationTheory.BoundStateRepresentation.Az}
  \uses{thm:hyd-rescaling-nonzero, thm:hyd-cross-linearity}
  $\vec A := \vec M/\sqrt{-2mE}$ on $V_E$, chosen so that $\vec J$ and
  $\vec A$ satisfy the commutation relations of two copies of $\liealg{so}_3$
  (Theorem~\ref{thm:hyd-two-factors-commute}).
\end{definition}

\section{The hidden $\sofour$ symmetry}
\label{sec:hyd-symmetry}

\begin{definition}[Two angular momentum triples]
  \label{def:hyd-two-triples}
  \lean{QuantumRepresentationTheory.BoundStateRepresentation.Ix, QuantumRepresentationTheory.BoundStateRepresentation.Iy, QuantumRepresentationTheory.BoundStateRepresentation.Iz, QuantumRepresentationTheory.BoundStateRepresentation.Kx, QuantumRepresentationTheory.BoundStateRepresentation.Ky, QuantumRepresentationTheory.BoundStateRepresentation.Kz}
  \uses{def:am-ladder, def:hyd-rescaled-rl}
  On $V_E$, set $\vec I := \tfrac12(\vec J + \vec A)$ and
  $\vec K := \tfrac12(\vec J - \vec A)$.
\end{definition}

\begin{theorem}[$I$ and $K$ each satisfy $\liealg{so}_3$ relations and commute with each other]
  \label{thm:hyd-two-factors-commute}
  \lean{QuantumRepresentationTheory.BoundStateRepresentation.I_comm_xy, QuantumRepresentationTheory.BoundStateRepresentation.I_comm_yz, QuantumRepresentationTheory.BoundStateRepresentation.I_comm_zx, QuantumRepresentationTheory.BoundStateRepresentation.K_comm_xy, QuantumRepresentationTheory.BoundStateRepresentation.K_comm_yz, QuantumRepresentationTheory.BoundStateRepresentation.K_comm_zx, QuantumRepresentationTheory.BoundStateRepresentation.I_comm_K}
  \uses{def:hyd-two-triples, thm:am-commutators}
  $[I_x,I_y]=iI_z$ (and cyclic permutations), likewise for $\vec K$, and
  $[I_a, K_b] = 0$ for all $a,b \in \{x,y,z\}$.
\end{theorem}

\begin{theorem}[$\sofour \cong \liealg{so}_3\oplus\liealg{so}_3$ decomposition]
  \label{thm:hyd-so4-decomposition}
  \lean{QuantumRepresentationTheory.BoundStateRepresentation.IRep, QuantumRepresentationTheory.BoundStateRepresentation.KRep}
  \uses{thm:hyd-two-factors-commute, thm:am-is-sl2-triple}
  $\vec I$ and $\vec K$ each generate an $\sltwo$-triple in the sense of
  Theorem~\ref{thm:am-is-sl2-triple}, realizing the dynamical symmetry
  algebra of the bound states as $\liealg{so}_3 \oplus \liealg{so}_3$.
\end{theorem}

\begin{theorem}[Equality of the two Casimir eigenvalues]
  \label{thm:hyd-equal-casimirs}
  \lean{QuantumRepresentationTheory.BoundStateRepresentation.equal_casimirs}
  \uses{thm:am-j2-casimir, def:hyd-two-triples}
  On any irreducible summand of $V_E$ under the joint $(\vec I, \vec K)$
  action, the Casimir eigenvalues of the $\vec I$-triple and the $\vec
  K$-triple coincide: $i(i+1) = k(k+1)$, forcing $i=k$.
\end{theorem}

\section{The energy spectrum}
\label{sec:hyd-spectrum}

\begin{theorem}[Classification of the bound-state representation]
  \label{thm:hyd-classification}
  \lean{QuantumRepresentationTheory.BoundStateRepresentation.classification_finrank}
  \uses{thm:sl2-commuting-classification, thm:hyd-so4-decomposition}
  Each irreducible summand of $V_E$ under the commuting pair
  $(\vec I,\vec K)$ is isomorphic to $V(2i)\otimes V(2k)$, by
  Theorem~\ref{thm:sl2-commuting-classification} applied to the two
  commuting $\sltwo$-triples of Theorem~\ref{thm:hyd-so4-decomposition}.
\end{theorem}

\begin{theorem}[Equal Casimirs force equal highest weights]
  \label{thm:hyd-equal-hw}
  \lean{QuantumRepresentationTheory.BoundStateRepresentation.dimension_eq_sq}
  \uses{thm:hyd-equal-casimirs, thm:hyd-classification}
  In the decomposition of Theorem~\ref{thm:hyd-classification}, the summand
  occurring in $V_E$ has $i = k$.
\end{theorem}

\begin{theorem}[Dimension of the bound-state eigenspace]
  \label{thm:hyd-dimension}
  \lean{QuantumRepresentationTheory.BoundStateRepresentation.dimension_eq_sq}
  \uses{thm:hyd-equal-hw, thm:sl2-basis-theorem}
  With $i=k$, the summand $V(2i)\otimes V(2i)$ has dimension $(2i+1)^2$; the
  quantity $n := 2i+1$ is a positive integer.
\end{theorem}

\begin{theorem}[Quadratic energy equation]
  \label{thm:hyd-quadratic-energy}
  \lean{QuantumRepresentationTheory.BoundStateRepresentation.quadratic_energy}
  \uses{thm:hyd-rescaling-nonzero, thm:hyd-dimension}
  Expressing the original Hamiltonian $H$ in terms of $\vec J^2 + \vec A^2$
  and the shared Casimir eigenvalue $i(i+1)$ from
  Theorem~\ref{thm:hyd-equal-hw} yields
  \[
    E = -\frac{me^4}{2\hbar^2 n^2}, \qquad n = 2i+1 \in \{1,2,3,\dots\}.
  \]
\end{theorem}

\begin{theorem}[Hydrogen energy levels and degeneracy]
  \label{thm:hyd-energy-degeneracy}
  \lean{QuantumRepresentationTheory.BoundStateRepresentation.energy_degeneracy}
  \uses{thm:hyd-quadratic-energy, thm:hyd-dimension}
  The bound-state energies of the hydrogen atom are
  $E_n = -me^4/(2\hbar^2 n^2)$ for $n \in \mathbb{Z}_{>0}$, and the
  eigenspace $V_{E_n}$ has dimension $n^2$.
\end{theorem}

\chapter*{Caveats}

This chapter is not part of the theorem dependency graph above; it collects
things found while auditing this blueprint against the current state of
Mathlib, before making it public. Some are corrections already applied
above; others are open opportunities for whoever picks this up next.

\paragraph{The Casimir element does not exist in Mathlib.} Checked with a
case-insensitive search of the entire Mathlib tree for ``Casimir'': zero
hits. Everything in Chapter~\ref{ch:casimir} (basis-dependent construction,
change-of-basis invariance, centrality via the universal enveloping algebra)
is new relative to Mathlib, not merely restated in a different form.

\paragraph{No exceptional low-rank Lie algebra isomorphisms.} Mathlib has a
genuine general construction of $\liealg{so}(n)$ as skew-symmetric matrices
(\texttt{LieAlgebra.Orthogonal.so}, with the equivalent ``indefinite form''
presentation \texttt{so'}, covering $\liealg{so}(3)$ and $\liealg{so}(4)$ as
the low-rank cases of types $B_1$ and $D_2$), and separately the general
skew-adjoint-endomorphism construction
(\texttt{Mathlib.Algebra.Lie.SkewAdjoint}) it is built from. But we found no
proof anywhere of the classical exceptional isomorphisms
$\liealg{so}(3)\cong\sltwo$ or $\liealg{so}(4)\cong\sltwo\oplus\sltwo$ ($B_1
\cong A_1$, $D_2\cong A_1\times A_1$) --- we checked the root-system and
Dynkin-diagram files
(\texttt{CartanMatrix.lean}, \texttt{Irreducible.lean}, even the dedicated
\texttt{RootSystem/Finite/G2.lean} for the $G_2$ exceptional type) and found
nothing for rank $\le 2$. Theorem~\ref{thm:am-is-sl2-triple} and the
$\vec I,\vec K$ construction in \S\ref{sec:hyd-symmetry} therefore prove
these isomorphisms \emph{ad hoc}, computationally, rather than by invoking a
general theorem, because no such general theorem is available to invoke.

\paragraph{Mathlib already has the concrete Lie ring structure this project
needs for $\liealg{so}(3)$, unused.} \texttt{Mathlib.LinearAlgebra.CrossProduct}
packages \texttt{Fin 3 → R} with the cross product as a genuine
\texttt{LieRing} (\texttt{Cross.lieRing}), with antisymmetry
(\texttt{cross\_anticomm}) and the Jacobi identity (\texttt{jacobi\_cross})
already proved, and the cross product itself is bilinear by construction
(it is built as a \texttt{LinearMap.mk₂}, so
Theorem~\ref{thm:hyd-cross-linearity} would be free). Neither
\texttt{AngularMomentumRepresentation} (\S\ref{sec:am-basic}) nor the
Levi-Civita-based \texttt{crossComponent} (\S\ref{sec:hyd-basic}) currently
uses this: angular momentum is presented as three operators satisfying
axiomatized commutation relations, rather than as a representation of this
concrete, already-available Lie ring. Routing Chapters~\ref{ch:am}--\ref{ch:hydrogen}
through \texttt{Cross.lieRing} instead is a live simplification opportunity,
not attempted in this pass.

\paragraph{Schur's lemma exists in Mathlib at two levels, neither directly
reusable yet.} \texttt{Mathlib.RingTheory.SimpleModule.Basic} gives the
division-ring level (\texttt{bijective\_or\_eq\_zero},
\texttt{Module.End.instDivisionRing}) but stops there --- no
algebraically-closed-field ``it's a scalar'' step. That step exists in full
generality in \texttt{Mathlib.CategoryTheory.Preadditive.Schur}
(\texttt{endomorphism\_simple\_eq\_smul\_id}), for any $\Bbbk$-linear
category with kernels, algebraically closed $\Bbbk$, finite-dimensional hom
spaces --- exactly Theorem~\ref{thm:rep-schur}'s statement, one level of
abstraction up. The obstruction to reusing it directly is that Mathlib has
no bundled category of Lie modules (no \texttt{LieModuleCat}, unlike
\texttt{ModuleCat} or \texttt{FGModuleCat}). Group representations do not
have this problem: \texttt{Mathlib.RepresentationTheory.FDRep} is a genuine
category and imports \texttt{CategoryTheory.Preadditive.Schur} directly to
get Schur's lemma for free (\texttt{finrank\_hom\_simple\_simple}).
\texttt{FDRep} is the exact template a \texttt{LieModuleCat} would follow;
building it is a well-scoped, plausible Mathlib contribution, but
Theorem~\ref{thm:rep-schur} as it stands takes the direct route (extract an
eigenvalue via \texttt{Module.End.exists\_eigenvalue}, show the shifted map
is a \texttt{LieModuleHom} with nontrivial kernel, conclude by
irreducibility) rather than waiting on that infrastructure.

\paragraph{No Clebsch--Gordan or weight-multiplicity machinery.} Checked for
any existing tensor-product decomposition or weight-multiplicity result for
$\sltwo$ or semisimple Lie algebra representations (searching for
``ClebschGordan'', and for weight-space/tensor-product combinations in
\texttt{Mathlib.Algebra.Lie.Weights}): nothing found. Everything in
\S\ref{sec:sl2-commuting}--\S\ref{sec:sl2-cg} is new.

\paragraph{The Killing form claim was wrong in an earlier draft, now fixed.}
An earlier version of Definition~\ref{def:cas-invariant-form}'s parenthetical
about the Killing form asserted a fact about Mathlib without checking it
closely enough, then a hasty correction asserted the opposite, also without
checking closely enough. The actual state, now reflected there: Mathlib has
both Cartan criteria in \texttt{Mathlib.Algebra.Lie.CartanCriterion}, and
\texttt{LieAlgebra.hasTrivialRadical\_iff\_isKilling} does give the full
equivalence between trivial radical (which coincides with semisimplicity in
characteristic zero) and a nondegenerate Killing form. Flagged here as a
reminder that citations in this blueprint should be checked against the
actual Mathlib source, not assumed from familiarity with the mathematics.

\paragraph{``Self-adjoint'' was removed from Definition~\ref{def:am-representation}.}
An earlier draft described the angular momentum operators as self-adjoint,
which does not match the actual Lean definition (deliberately: physical
self-adjointness is exactly the kind of fact that belongs to a future
Physlib bridge, not to this algebraic model). The blueprint text is now
corrected to match the Lean code.

\paragraph{The explicit-representation bridge (\S\ref{sec:rep-basic},
Remark~\ref{def:rep-bridge}) remains unbuilt.} This is the one gap flagged
already in Chapter~1 that is worth re-emphasizing here: Chapters~4--5 use
\texttt{LieRepresentation} (an explicit morphism into $\End_K(V)$) rather
than Mathlib's instance-based \texttt{LieModule}, specifically because
Mathlib has no public constructor turning an explicit morphism back into
local \texttt{LieRingModule}/\texttt{LieModule} instances. Until that bridge
is built, results proved in the generic, instance-based Chapters~1--3
(irreducibility transfer, the classification theorems, Casimir scalars) have
to be re-derived or manually transported for the explicit-morphism objects
in Chapters~4--5 rather than applied directly.
